\section{Experimental Setup}
\label{sec:setup}

\subsection{Artifacts, and what each may be used for}
\label{sec:artifacts}

The evidence in this paper comes from five committed artifacts that are \emph{not}
interchangeable. Table~\ref{tab:datasets} lists them; we describe the
distinctions that matter, because conflating them is the single easiest way to
produce a wrong number.

% GENERATED by scripts/paper/make_tables.py -- do not edit by hand.
\begin{table}[t]
\centering
\small
\caption{Committed experiment artifacts. Behavioral rows never enter an official-success denominator.}
\label{tab:datasets}
\begin{adjustbox}{max width=\linewidth}
\begin{tabular}{lrrl}
\toprule
Dataset & Records & Design & Evaluation \\
\midrule
Frozen ReAct control baseline & 48 & 16 tasks $\times$ 3 seeds & official; 25 success \\
Behavioral fault matrix & 800 & 5 faults $\times$ 16 tasks $\times$ 5 seeds $\times$ 2 & no HAR; behavioral only \\
Official fault subset & 400 & 5 faults $\times$ 8 tasks $\times$ 5 seeds $\times$ 2 & official; 256 success \\
Architecture pilot & 54 & 3 tasks $\times$ 3 seeds $\times$ 3 conditions $\times$ 2 arch & official; PILOT \\
Recovery subset & 45 & 5 faults $\times$ 3 tasks $\times$ 3 seeds & log heuristics; subset of the official 400 \\
\bottomrule
\end{tabular}
\end{adjustbox}
\end{table}


Three distinctions carry real weight. The \textbf{official fault subset}
(\OfficialN{} evaluations) is the only source for success-rate claims, and it
contains no step or timing columns. The \textbf{behavioural fault matrix}
(\TotalTrials{} records) retained no HAR, so it carries no official verdict at
all --- every row is marked \texttt{NOT\_EVALUATED\_NO\_HAR} --- and it supports
process-measure claims and nothing else; note also that its \TotalTrials{}
records are \MatrixPairs{} pairs logged twice, not independent samples. The
\textbf{official control baseline} (\BaselineN{} trials) is the only
clean-condition artifact carrying both an official verdict and process measures.
The \textbf{architecture pilot} (\ArchN{} rows) and \textbf{recovery subset}
(\RecN{} rows) are small exploratory artifacts, treated as pilots throughout.

Two consequences we enforce mechanically rather than by convention. First,
official success is never aggregated across artifacts: the behavioural matrix
cannot enter a success-rate denominator, which is why the process results and the
success results below are reported from different data and never combined into a
single sentence about ``cost per unit of success.'' Second, the baseline and the
official subset share a column schema but not a design (\BaselineTasks{} tasks
$\times$ \BaselineSeeds{} seeds versus \OfficialTasks{} $\times$ \TrialsPerCell{}),
so rows are tagged with their source artifact and the two are never concatenated.

\subsection{The evaluator has two paths, and they are not equivalent}
\label{sec:fallback}

The official evaluator in this stack has a native path and a compatibility
fallback used when the native retrieval check cannot be applied. The fallback is
\emph{structurally} failure-biased: its null-retrieval schema branch returns
success only when the state matches \emph{and} both retrieved values are null,
which our runs almost never satisfy. Empirically, of the \OfficialFallbackN{}
evaluations it processed in the official subset, \textbf{zero} succeeded, and the
same pattern holds in the baseline (\BaseFallbackRate{} against
\BaseNativeRate{} on the native path).

We therefore stratify every official result by evaluator path and treat the
native stratum as primary. Table~\ref{tab:baseline} shows the effect on the
frozen baseline: the pooled rate is \BasePooledRate{}, the native rate is
\BaseNativeRate{}, a \BaselineSwing{}~pp gap produced entirely by which trials
happened to be scored by which path. Reported without stratification, that gap
would be attributed to the agent.

Because each arm of a pair may be scored by a different path, the native
analysis restricts to pairs where \emph{both} arms were scored natively. This is
why the per-fault $n$ drops from \TrialsPerCell{} seeds $\times$
\OfficialTasks{} tasks $= 40$ to \ResWebDomMissingNNat--\ResWebTimeoutNNat{}. We
report both the pair-restricted rates (the quantity $\Delta$ and McNemar refer
to) and the per-arm marginal native rates, and we label them separately in
Table~\ref{tab:official_native} because they are computed over different subsets
and are not interchangeable.

\subsection{Metrics}
\label{sec:metrics}

\emph{Success rate} is the official evaluator's verdict. \emph{Paired risk
difference} $\Delta$ is in percentage points, fault minus control, on the
pair-restricted subset. \emph{Process cost} is the within-pair delta in
\texttt{step\_count} and wall-clock seconds. \emph{Discordance} $\pi_d$ is the
fraction of pairs in which the two arms disagree; it is the quantity that governs
power, and it is reported rather than assumed.

\subsection{Reproducibility of the runs}
\label{sec:runs}

Every trial records its fault seed, injection step, model profile, architecture,
and git commit, so a fault condition is reproducible at the level of the design
rather than the token sequence --- the provider does not guarantee bit-identical
sampling for a fixed seed. Wall-clock is host- and load-dependent, which is why
process-cost claims use within-pair deltas only; \appref{app:runs} quantifies the
drift.
